% Chapter 5. Numbers only with status ✅ in docs/research/traceability.md and a passport in
% docs/thesis/figures.md. Pilot results are marked \pending until they exist.
% Preliminary checks: docs/research/research-log.md (2026-10-17), docs/requirements/quality-scenarios.md,
% docs/architecture/change-scenarios.md. RQ-4 incidents: docs/thesis/materials-log.md, docs/journal/sessions/.
\chapter{Evaluation and Discussion}\label{ch:eval}

This chapter reports what has been measured so far and what the pilot will add. Section~\ref{sec:eval-status}
summarises the state of each study. Section~\ref{sec:eval-prelim} presents preliminary technical
validation. Sections~\ref{sec:eval-rq1}--\ref{sec:eval-rq3} describe the results for RQ-1 to RQ-3, which
depend on pilot data, and Section~\ref{sec:eval-rq4} the evidence collected so far for RQ-4.
Section~\ref{sec:eval-discussion} discusses the findings and Section~\ref{sec:eval-limits} their limits.

\section{Status of the evaluation}\label{sec:eval-status}

\begin{table}[ht]
\centering\small
\caption{State of the four studies.}
\label{tab:eval-status}
\begin{tabularx}{\linewidth}{@{}P{2.6cm}LL@{}}
\toprule
Study & Ready now & Needs \\
\midrule
1. Baseline (RQ-1) & Annotation in Telegram and packages; phoneme-level metrics, $\kappa$, $\alpha$
computed automatically; error catalogue; confusion and book-line reports; entry questionnaire & Pilot
attempts; two annotators; bootstrap intervals; confirmed-only frequencies \\
2. Adaptation (RQ-2) & Raw model output and versions in every attempt; recogniser port; phoneme-level
metrics & Training pipeline (TD-FN-72); blind test annotation (TD-FN-73); GPU (TD-FN-74); pilot data \\
3. Feedback (RQ-3) & Decision trace; retry data; preregistered model & Condition switch per card; delayed
test; pilot \\
4. Harness (RQ-4) & Hook journal (since 18 October 2026); session and prompt logs; ADRs; incidents &
Counts over the whole project; coding of incidents \\
\bottomrule
\end{tabularx}
\end{table}

\section{Preliminary technical validation}\label{sec:eval-prelim}

\subsection{Does the pipeline separate right from wrong reading?}

Before any learner used the tutor, the pipeline was checked on expert audio: the teaching recitation of
Mahmoud Khalil al-Husary \cite{SRC-042}. Four verses were sent as attempts for their own text, and three
verses as attempts for a different text. Table~\ref{tab:husary} shows the scores.

\begin{table}[ht]
\centering\small
\caption{Scores on expert recitation (17 October 2026, real model).}
\label{tab:husary}
\begin{tabular}{@{}llr@{}}
\toprule
Condition & Verses & Score \\
\midrule
Audio matches the text & 112:1 & 100 \\
 & 112:2 & 100 \\
 & 1:2 & 99 \\
 & 105:1 & 99 \\
\midrule
Audio of another verse & 112:1 vs.\ 112:2 & 67 \\
 & 1:2 vs.\ 112:1 & 40 \\
 & 105:1 vs.\ 1:2 & 48 \\
\bottomrule
\end{tabular}
\end{table}

Correct readings scored 99--100 and mismatched texts 40--67, well apart from the pass threshold of 80. This
shows that the pipeline --- conversion, recognition, phonetisation, alignment and grading --- works end to
end and discriminates correct from incorrect reading on expert audio. It says nothing yet about beginners:
the speaker is an expert, the text is Qur'anic, and the sample has seven pairs. A synthetic Arabic voice
of the operating system was recognised poorly, so synthetic speech cannot replace human recordings in tests.

\subsection{Response time}

On the first nine attempts with the real model on the author's laptop, the median processing time was
1.7\,s and the 90th percentile 5.0\,s, within the 10-second target of QAS-03. The measured interval covers
recognition and assessment only; downloading the voice message, conversion and rendering the answer are
not included, so the learner waits longer than these numbers. The pilot will report the same percentiles per
day on the deployment server.

\subsection{Robustness of the assessment algorithm}

The model can return any sequence, so the assessment must behave sensibly for all of them. The six
properties listed in Section~\ref{sec:impl-verification} were checked on random phoneme strings: 300 cases
per property in every run of the suite and 5\,000 per property in a one-off run, without a violation.

\subsection{Modifiability}

Adding a second recogniser required three files and no change to the domain, the use cases, the database
or the interface; all architecture contracts and tests passed (QAS-10). This supports the claim that the
tutor can be used to compare models on the same recordings; this is exactly how RQ-2 compares B0, B1 and
B2.

\section{RQ-1: the baseline model and the error profile}\label{sec:eval-rq1}

\begin{notebox}
This section will be completed after the pilot and double annotation. The tables below show the planned
layout; no value is reported before it has a passport (report, date, $n$, code fingerprint).
\end{notebox}

\begin{table}[ht]
\centering\small
\caption{Quality of error detection on beginners (T-01).}
\label{tab:T-01}
\begin{tabular}{@{}lccc@{}}
\toprule
Measure & Value & 95\,\% CI & $n$ \\
\midrule
Precision & \pending{T-01} & & \\
Recall & & & \\
F1 & & & \\
Diagnosis accuracy & & & \\
FAR (verdict) & & & \\
FRR (verdict) & & & \\
\bottomrule
\end{tabular}
\end{table}

\begin{table}[ht]
\centering\small
\caption{Agreement on verdicts (T-02).}
\label{tab:T-02}
\begin{tabular}{@{}lcc@{}}
\toprule
Pair & Cohen's $\kappa$ & 95\,\% CI \\
\midrule
Tutor -- teacher A & \pending{T-02} & \\
Tutor -- teacher B & & \\
Teacher A -- teacher B & & \\
\bottomrule
\end{tabular}
\end{table}

The analysis will report precision and recall by error type and by letter (R-08), Krippendorff's $\alpha$ per
phoneme between annotators, the share of disputed attempts that annotators corrected, and the
self-correction block separately. H1.1--H1.3 will be tested as preregistered, H1.3 also on the blindly
annotated part.

\paragraph{Error profile} \pending{error-profile results}. The analysis will present the frequency of
error kinds (T-03), the most difficult lines and minimal pairs of the book (T-04), and confusion matrices
for all reported and for confirmed errors (R-09). The comparison of the two matrices shows which parts of
the apparent learner profile are produced by the model, and also suggests on which sounds adaptation may
gain most. H1.4 needs the attribute outputs of the model, which the tutor already receives but does not yet
store separately; the comparison by first language (E-1) is run only if both groups are present.

\section{RQ-2: model adaptation}\label{sec:eval-rq2}

\begin{notebox}
This section will be completed after B1 and B2 are trained and the test set is run once (April 2027).
Below is the planned layout of the tables; values will appear only with a passport.
\end{notebox}

\begin{table}[ht]
\centering\small
\caption{B0, B1 and B2 on test learners with blind gold labels (T-11).}
\label{tab:T-11}
\begin{tabular}{@{}lccc@{}}
\toprule
Measure & B0 & B1 & B2 \\
\midrule
Precision / recall / F1 & \pending{T-11} & & \\
FAR / FRR (verdict) & & & \\
$\kappa$ with gold labels & & & \\
Share of gap closed $G_B$ & 0 & & \\
ECE & & & \\
Share of ``not sure'' & --- & & \\
\bottomrule
\end{tabular}
\end{table}

\begin{table}[ht]
\centering\small
\caption{Retention on expert recitation: PER of B0 and B2 (T-12, H2.3).}
\label{tab:T-12}
\begin{tabular}{@{}lccc@{}}
\toprule
Set & B0 & B2 & Difference, 95\,\% CI \\
\midrule
Held-out expert recitation & \pending{T-12} & & \\
\bottomrule
\end{tabular}
\end{table}

The analysis will report the differences B1 $-$ B0 and B2 $-$ B0 with paired cluster-bootstrap intervals
(H2.1, H2.2), the non-inferiority check on expert recitation (H2.3), whether the conformal guarantee holds
on the test set, B1's risk--coverage curve (R-14) and B2's learning curve (R-15). If the fallback of
Section~\ref{sec:method-s2} applies, B2 will be described as a feasibility check and RQ-2 will be answered
by B1.

\section{RQ-3: effect of explanations}\label{sec:eval-rq3}

\pending{RQ-3 results}. Before the pilot the tutor needs one more feature: a per-learner, per-card switch
between the full explanation and the correct/incorrect signal, with the condition recorded in the attempt.
The analysis will report the score change after each condition (R-11), the recurrence of named errors and the
delayed-test results (T-05), with the mixed model of Section~\ref{sec:method-s3}. Without the control
condition, an improvement on retry would mix the effect of the explanation with the effect of simply trying
again; for this reason descriptive retry statistics collected before the switch will not be used as evidence.

\section{RQ-4: the development harness}\label{sec:eval-rq4}

Quantitative measures --- blocks by reason (R-13), interventions per session and iterations until green
(T-08) --- will be computed over the whole project at its end; the hook journal started on 18 October 2026.
\pending{RQ-4 counts}. The journals already contain incidents that show what the harness caught, what it
missed and where the human was needed. Table~\ref{tab:incidents} groups them by failure mode.

\begin{table}[ht]
\centering\small
\caption{Documented incidents from the session journals (October 2026).}
\label{tab:incidents}
\begin{tabularx}{\linewidth}{@{}P{3.0cm}LP{2.6cm}@{}}
\toprule
Failure mode & Incident & Caught by \\
\midrule
Test gaming & An accepted test contained a wrong expectation written by the agent; the change was made
only through the tests phase with the author's consent. & Frozen tests, stop signal \\
Architecture erosion & A new research module imported the bootstrap layer; the entry point was moved. &
import-linter \\
Leak of secrets & A regular expression written by the agent did not mask the bot token inside a URL. & Test
on logs \\
Silent defect & Dates lost their time zone in SQLite; a UTC-aware column type was added. & Integration
test \\
Destructive edit & While removing a keyboard, the agent truncated the file with the consent keyboards. &
Type checker \\
Unspoken assumption & The bot was handed to the author with the stub recogniser, which accepts everything;
the author saw ``all correct''. & Not caught; a hook now reminds the agent \\
Rule left as prose & Consent button texts lay outside the protected texts file, so the agent could change
them without approval. & Review; moved and tested \\
Check with a blind spot & An approval rule covered editing but not writing the same file. & Review \\
Check with a blind spot & A path rule in \texttt{.gitignore} hid five documentation files from every
commit. & Consistency audit \\
Check with a blind spot & The dependency audit did not see the optional model extra and its
vulnerabilities. & Dependency review \\
Green but broken & The container could not write logs, and a CI job lacked a test dependency; both passed
all local checks. & Direction reviewer (STOP) \\
Text does not match action & A button promised the alphabet while the lesson started with the
\emph{basmala}. & Console bot driver \\
\bottomrule
\end{tabularx}
\end{table}

Three observations already stand out. First, the deterministic checks caught the classic agent failures:
test gaming was blocked and handled through the author, architecture violations and type errors stopped
the agent. Second, several checks were green while blind: a configuration rule, an ignore pattern and a
dependency audit each covered less than they appeared to. Each was found by another check or a review, which
argues for checking the checks. Third, the most costly incident --- the stub recogniser --- was a failure to
communicate, not to compute; it was fixed by adding a hook that turns the reminder into a mechanism. The
quantitative analysis will show how often such incidents occur relative to the volume of work.

\section{Discussion}\label{sec:eval-discussion}

\paragraph{Two probabilistic systems under one discipline} The thesis deals with two systems that can pass a
test without doing the right thing. The coding agent can make tests pass by changing them; the recogniser can
reject a correct reading or accept a wrong one, and an adapted model can show a gain that comes from
leakage rather than learning. The remedies have the same shape: deterministic external
checks that the system cannot edit, a log of every decision, and an escalation path to a human. For the
agent these are frozen tests, layer contracts and STOP verdicts; for the model they are the decision trace,
the ``the bot was wrong'' button, teacher annotation, a frozen and blindly annotated test set, and a
``not sure'' outcome with a guarantee on the false-rejection rate.

\paragraph{What a tutor can and cannot claim} The literature and the preliminary checks support a modest
role: a patient practice partner that gives immediate, explicit first-line feedback between lessons. The
claim that it replaces a teacher is neither supported by evidence nor consistent with the book. The design
reflects this by making disputes easy, sending disputed attempts to teachers first and recording everything
a teacher needs to review a verdict.

\paragraph{Pending findings} \pending{interpretation of RQ-1--RQ-3 results, including how much of the gap to a
teacher adaptation closed and how much data it needed}.

\section{Limitations}\label{sec:eval-limits}

The pilot is small and closed, so its results describe the participants rather than all Russian-speaking
beginners. Annotators of the training part see the tutor's answer, which may inflate agreement; the test
set is therefore annotated blind. The test set is small --- about ten learners --- so intervals for RQ-2 will
be wide, and the gain from adaptation on these learners need not transfer to other Russian speakers or to
later lessons of the book. The model was trained on Qur'anic
recitation, whereas pages 10--37 of the book contain syllables and words that are not Qur'anic; results are
therefore reported by material type. The similarity score used for verdicts is coarse on long lines. The
case study of the harness has one developer, one project and an author who is also the evaluator; it is
illustrative, not comparative. Finally, the curriculum text transcribed from the scan is still awaiting full
human verification; every attempt stores the curriculum version, so attempts affected by a later correction
can be identified and excluded.
